% META: {"id":"TCOMPL-MF-EX-003","chapitre":"TCOMPL-MODELES-FONCTION","type_objet":"exercice","capacites":["TCOMPL-MF-C5"],"capacites_codes":["C5"],"methodes":[],"parcours":2,"competences":["calculer"],"duree_min":15,"mode_creation":"ex_nihilo","sources_inspiration":[],"fichier_tex":"chapitres/TCOMPL-MODELES-FONCTION/exercices/TCOMPL-MF-EX-003.tex","corrige_tex":"chapitres/TCOMPL-MODELES-FONCTION/corriges/TCOMPL-MF-CO-003.tex","status":"approved","origin":"generated","status_history":["generated","audited","accepted","approved"],"release_acceptance":"RELEASE_OWNER_BATCH_ACCEPTANCE","acceptance_closure_digest":"sha256:9b3ccf9a81c5520fbb7e03b7b2d3e7bf2057a02834b6d908bf3be5f49f3b553f","acceptance_receipt_digest":"sha256:4fad86f0282e0fa4e0419cca1ad6379ae7af5a280c04a4a90880f90a1c044242"}
% BEGIN-VERIFY
% from sympy import symbols, diff, solve, Eq, simplify
% x = symbols('x')
% h = x**3 - 6*x**2 + 9*x
% hpp = diff(h, x, 2)
% assert simplify(hpp - (6*x - 12)) == 0
% assert solve(Eq(hpp, 0), x) == [2]
% assert hpp.subs(x, 0) < 0
% assert hpp.subs(x, 3) > 0
% END-VERIFY
\begin{exercice}{TCOMPL-MF-EX-003}{2}{15}
Soit $h(x) = x^3 - 6x^2 + 9x$ sur $\mathbb{R}$.
\begin{enumerate}
  \item Calculer $h''(x)$.
  \item Étudier le signe de $h''(x)$, et en déduire les intervalles où $h$ est convexe, où $h$ est concave.
  \item $h$ admet-elle un point d'inflexion ? Si oui, préciser son abscisse.
\end{enumerate}
\end{exercice}
